% !TEX root = hw2-answers.tex
\chead{\textbf{Exercises week 2}}

\begin{document}
\flushleft\textbf{Topic: Markov kernels}

\begin{enumerate}
\item ($6 \times 1/3$ points) Consider Player 1 and Player 2, playing multiple games of chess.
Let $\Tcal = \Wcal = \Zcal = \{1,2\}$ and let $T, W$ and $Z$ denote the winning player of the first, second and third game respectively. When a player wins a game, the other player gets to play first (white) in the following game, providing a significant advantage over the other player (black). A coin toss represented by $C \in \Ccal = \{1,2\}$, decides which player gets to play white in the first game.

Provide an interpretation of the following Markov kernels: % Add next year (2023): Provide, in words, an interpretation of ...
\begin{enumerate}
  \item the Markov kernels $N(Z|W), M(W|T), Q(T|C)$ and $K(C|*)$
  \item the composition $N(Z|W) \circ \big(M(W|T)\circ Q(T|C)\big)$
  \item the product $P(Z, W, T|C) := N(Z|W) \otimes \big(M(W|T)\otimes Q(T|C)\big)$
  \item the marginal $P(W|C)$; express this as a composition of two Markov kernels
  \item the conditional Markov kernel $P(Z,T|W,C)$, and how does this relate to the product and marginal in (c) and (d)?
\end{enumerate}
What the player's don't know, is that after three games the player with the highest number of wins receives a prize. However, from the prize money a fine is subtracted when the winning player has won the coin flip and thus played white in the first game. Accordingly, let
$$f(z,w,t, c) = \begin{cases} 100 - 10*\mathbbm{1}_{\{1\}}(c) & \text{if $z+w+t \leq 4$} \\ 0 &\text{otherwise}\end{cases}$$
denote the money won by Player 1. Using this, also provide an interpretation of the following quantity:
\begin{enumerate}
  \item[(f)] the push-forward $f_*P(Z, W, T|C)$
\end{enumerate}

\begin{gradedsolution}
  \begin{enumerate}
    \item
    \begin{enumerate}
      \item $N(Z|W)$: The distribution over who wins the third game, given the winner of the second game.
      \item $M(W|T)$: The distribution over who wins the second game, given the winner of the first game.
      \item $Q(T|C)$: The distribution over who wins the first game, given who got to play first, as decided by the outcome of the coin flip.
      \item $K(C|*)$: The distribution over the possible outcomes of the coin flip.
    \end{enumerate}
    \item The distribution over who wins the third game, given the outcome of the coin flip.
    \item The joint distribution over the winners of the first, second and third game, given the outcome of the coin flip.
    \item The distribution over who wins the second game given the outcome of the coin flip. It is equal to the composition $M(W|T)\circ Q(T|C)$.
    \item The joint distribution over wins the first and third game, given the outcome of the second game and the coin flip. It related to (c) and (d) via $P(Z, W, T|C) = P(Z, T|W, C) \otimes P(W|C)$
    \item The distribution over the prize money won by Player 1, given the outcome of the coin flip. Thus, when $C=1$, this is a distribution over $\{90,0\}$, and when $C=2$ this is a distribution over $\{100,0\}$.
  \end{enumerate}
\end{gradedsolution}

\item ($4 \times 1$ points) Let $\Tcal$ be a measurable space, let $\Xcal, \Ycal$ be standard measurable spaces, and consider the measurable functions 
\begin{align*}
  f_1&: \Wcal \rightarrow \Tcal \\
  f_2&: \Tcal \rightarrow \Xcal \\
  f_3&: \Tcal \rightarrow \Ycal
\end{align*}
and their Markov kernel representations $\delta_{f_1}, \delta_{f_2}, \delta_{f_3}$. For each of the following deterministic Markov kernels, specify its domain and codomain, express the function ($f_c/f_p/f_m/f_{cd}$) in terms of $f_1, f_2$ and $f_3$, and show why this is correct:
% For in 2023: “express the functions $f_c/f_p/f_m/f_{cd}$ in terms of $f_1, f_2$ and $f_3$”


\begin{enumerate}
  \item the composition $\delta_{f_c} := \delta_{f_2} \circ \delta_{f_1}$
  \item the product $\delta_{f_p} := \delta_{f_2} \otimes \delta_{f_3}$
  \item $\delta_{f_m}$ as the $X$-marginal of $\delta_{f_p}$
  \item the conditional $\delta_{f_{cd}}$ such that $\delta_{f_p} = \delta_{f_{cd}} \otimes \delta_{f_m}$
\end{enumerate}

\begin{gradedsolution}
  \begin{enumerate}
    \item From Remark \ref*{deterministic-mk-composition} it directly follows that $f_c = f_2 \circ f_3$. More specifically, we have
    $$\delta_{f_2}(X|T)\circ \delta_{f_1}(T|W) : \Wcal \dshto \Xcal,$$ and for all $w\in \Wcal$ and $C\in \Bcal_{\Xcal}$ we have
    \begin{align*}
      [\delta_{f_2}(X|T)\circ \delta_{f_1}(T|W)](C, w) &= \int_\Tcal\int_\Xcal \mathbbm{1}_C(x)\delta_{f_2}(dx|T=t)\delta_{f_1}(dt|W=w) \\
      &= \int_\Tcal \delta_{f_2}(C|T=t)\delta_{f_1}(dt|W=w) \\
      &= \int_\Tcal \mathbbm{1}_C(f_2(t))\delta_{f_1}(dt|W=w) \\
      &= \int_\Tcal \mathbbm{1}_{f_2^{-1}(C)}(t)\delta_{f_1}(dt|W=w) \\
      &= \delta_{f_1}(T\in f_2^{-1}(C)|W=w) \\
      &= \mathbbm{1}_{f_2^{-1}(C)}(f_1(w))) \\
      &= \mathbbm{1}_C(f_2(f_1(w))) \\
      &= \delta_{f_2\circ f_1}(X\in C|W=w) \\
      &= \delta_{f_2\circ f_1}(X|W)(C, w),
    \end{align*}
    so $f_c = f_2\circ f_1$.
    \item We use the following results:
    \begin{enumerate}
      \item 
      For standard Borel spaces $\Xcal$ and $\Ycal$ we have $\Bcal_{\Xcal\times \Ycal} = \Bcal_\Xcal \otimes \Bcal_\Ycal$.
      \item In general, the product sigma algebra $\Bcal_\Xcal \otimes \Bcal_\Ycal$ is generated by $\{A\times B: A\in\Bcal_\Xcal , B\in \Bcal_\Ycal\}$.
      \item If two measures $\mu_1, \mu_2$ agree on a $\pi$-system $\Ical$ (i.e.\ that $\mu_1(I) = \mu_2(I)$ for all $I\in\Ical$), then $\mu_1(B)=\mu_2(B)$ for all $B\in\sigma(\Ical)$.
      \item $\{A\times B: A\in\Bcal_\Xcal , B\in \Bcal_\Ycal\}$ is a $\pi$-system.
    \end{enumerate}
    So, we have to test the provided equality only on sets $A\times B: A\in\Bcal_\Xcal, B\in\Bcal_\Ycal$. For all such $A$ and $B$ and $t\in\Tcal$ we have
    \begin{align*}
      [\delta_{f_2}(X|T)\otimes \delta_{f_3}(Y|T)](A\times B, t) &= \int_\Xcal\int_\Ycal \mathbbm{1}_{A\times B}(x, y)\delta_{f_2}(dx|T=t)\delta_{f_3}(dx|T=t) \\
      &= \int_\Xcal\int_\Ycal \mathbbm{1}_{A}(x)\mathbbm{1}_{B}(y)\delta_{f_2}(dx|T=t)\delta_{f_3}(dx|T=t) \\
      &= \int_\Xcal \mathbbm{1}_{A}(x)\delta_{f_2}(dx|T=t)\int_\Ycal\mathbbm{1}_{B}(y)\delta_{f_3}(dx|T=t) \\
      &= \delta_{f_2}(A|T=t)\delta_{f_3}(B|T=t) \\
      &= \mathbbm{1}_{A}(f_2(t))\mathbbm{1}_{B}(f_3(t)) \\
      &= \mathbbm{1}_{A\times B}(f_2(t), f_3(t)) \\
      &= \delta_{(f_2, f_3)}(A\times B|T=t) \\
      &= [\delta_{(f_2, f_3)}(X, Y|T)](A\times B, t) \\
      &=: [\delta_{f_p}(X, Y|T)](A\times B, t),
    \end{align*}
    so $f_p = (f_2, f_3)$.
    \item For all $A\in \Bcal_\Xcal$ (and using an expression for $\delta_{f_p}$ from the previous exercise) we have
    \begin{align*}
      \delta_{f_m}(X|T)(A, t) &= \delta_{f_p}(X\in A, Y\in\Ycal|T=t) \\
      &= \int_\Xcal\mathbbm{1}_A(x)\delta_{f_2}(dx|T=t)\int_\Ycal\mathbbm{1}_\Ycal(y)\delta_{f_3}(dy|T=t) \\
      &= \int_\Xcal\mathbbm{1}_A(x)\delta_{f_2}(dx|T=t) \\
      &= \delta_{f_2}(A|T=t) \\
      &= \delta_{f_2}(X|T)(A, t),
    \end{align*}
    so $f_m = f_2$.
    \item Let $\delta_{f_{cd}}(Y|X, T): \Xcal\times \Tcal \dshto \Ycal$ be the conditional Markov kernel such that
    \begin{equation}\label{foo}
      \delta_{f_p}(X, Y|T) = \delta_{f_{cd}}(Y|X, T) \otimes \delta_{f_m}(X|T).
    \end{equation}
    Then, for all $A\in \Bcal_\Xcal$, $B\in\Bcal_\Ycal$ and $t\in\Tcal$ we have
    \begin{align*}
      [\delta_{f_3}(Y|T)\otimes \delta_{f_2}(X|T)](B\times A, t) &= [\delta_{f_2}(X|T)\otimes \delta_{f_3}(Y|T)](A\times B, t) &\text{Fubini} \\
      &= \delta_{f_p}(X, Y|T)(A\times B, t) &\text{(b)} \\
      &= [\delta_{f_{cd}}(Y|X, T) \otimes \delta_{f_m}(X|T)](B\times A, t) &\text{eqn.\ (\ref*{foo})} \\
      &= [\delta_{f_{cd}}(Y|X, T) \otimes \delta_{f_2}(X|T)](B\times A, t) &\text{(c)}
    \end{align*}
    and so for $f_{cd} = f_3$, by Lemma \ref*{ess-unique} we have that $\delta_{f_{cd}} = \delta_{f_3}$ up to $\delta_{f_m}(X|T)$-null sets.

  \end{enumerate}
\end{gradedsolution}

% \item (1 point) $\Zcal, \Tcal, \Wcal$ are assumed to be finite discrete spaces with Markov kernels:\\
 
%   \[Q(Z|W):\, \Wcal \dshto \Zcal,\quad (C,w) \mapsto Q(Z \in C|W=w),  \]
%   \[K(W|T):\, \Tcal \dshto \Wcal,\quad (D,t) \mapsto K(W \in D|T=t),  \]

% and corresponding stochastic matrices:

% \[\tilde Q = 
% \begin{bmatrix}
% 0.1 & 0.7 & 0.2 \\
% 0.3 & 0   & 0.4 \\
% 0.6 & 0.3 & 0.4 \\
% \end{bmatrix} 
% \qquad
% \tilde K = 
% \begin{bmatrix}
% 0.2 & 0.1 & 0.6 \\
% 0.3 & 0.1 & 0 \\
% 0.5 & 0.8 & 0.4 \\
% \end{bmatrix}\]

    
% Compute the stochastic matrix $\tilde P$ for the composed  Markov kernel: $P(Z|T):=Q(Z|W) \circ K(W|T)$.

% \begin{gradedsolution}
% $\tilde P = \begin{bmatrix}
% 0.33 & 0.24 & 0.14 \\
% 0.26 & 0.35 & 0.34 \\
% 0.41 & 0.41 & 0.52 \\
% \end{bmatrix}$.
% \end{gradedsolution}

\item (2 + 1 + 1 points) Consider two Markov kernels:
\[K(W|T):\, \Tcal \dshto \Wcal,\quad (D,t) \mapsto K(W \in D|T=t),  \]
\[Q(Z|W):\, \Wcal \dshto \Zcal,\quad (C,w) \mapsto Q(Z \in C|W=w),  \]
that have linear multivariate Gaussian densities (w.r.t.\ Lebesgue measure):
\begin{align*} 
  k(w|t) &= \Ncal(w|A \cdot t +b, \Sigma),  \\
  q(z|w) & = \Ncal(z| C \cdot w + d, \Gamma).
\end{align*}

Compute the parameters of the following multivariate Gaussian Markov kernels:
\begin{enumerate}
  \item the product $P(Z, W|T):=Q(Z|W) \otimes K(W|T)$
  \item the composition $P(Z|T):=Q(Z|W) \circ K(W|T)$
  \item the conditional $P(W|Z, T)$ of $P(Z, W|T)$
\end{enumerate}

Suggested steps:
\begin{enumerate}[label=\roman*)]
  \item First, recall how the general formula for the density of a composition is computed.
  \item Write the log joint density $\log p(z, w|t)$ as a quadratic polynomial in $\R^{\dim \Wcal + \dim \Zcal}$ using joint precision matrix $\Lambda$:
  \[ 
    -\frac{1}{2} \mat{z \\ w}^T  \mat{\Lambda_{zz} & \Lambda_{zw} \\ \Lambda_{zw}^T & \Lambda_{ww}}\mat{z \\ w} + \mat{b_1 \\ b_2}^T\mat{z \\ w} + \text{const}
  \]
  \item Invert to find the covariance matrices:
  \[ 
    -\frac{1}{2} \mat{z \\ w}^T  \mat{\Sigma_{zz} & \Sigma_{zw} \\ \Sigma_{zw}^T & \Sigma_{ww}}^{-1}\mat{z \\ w} + \mat{b_1 \\ b_2}^T\mat{z \\ w} + \text{const}
  \]
  \emph{Hint: You can use the following identity for symmetric matrices $M,N$ and matrix $T$ without proof:}
  \begin{align*}
    \mat{N &-NT \\ -T^TN  &M+T^TNT}^{-1} = \mat{N^{-1}+TM^{-1}T^T&  TM^{-1}  \\M^{-1}T^T & M^{-1} }
  \end{align*}
  \item Complete the square to obtain Gaussian form:
  \[ 
    \mat{z-\mu_z \\ w-\mu_w}^T \mat{\Sigma_{zz} & \Sigma_{zw} \\ \Sigma_{zw}^T & \Sigma_{ww}}^{-1}\mat{z-\mu_z \\ w-\mu_w} + \text{const}
  \]

  \item Read off the marginal $p(z|t)$.
  
  \item Find from external sources a way to derive the mean and covariance matrix for the conditional $P(W|Z, T)$. Clearly state what information you use from which external source, and how this translates to our setting.
\end{enumerate}

\begin{gradedsolution}
  \begin{enumerate}
    \item Define $\mu=At+b$.
    \begin{align*}
      &(z-Cw-d)^T \Gamma^{-1}(z-Cw-d)+(w-\mu)^T\Sigma^{-1}(w-\mu) \\
      &=\mat{z \\w}^T \mat{\Gamma^{-1} & -\Gamma^{-1}C \\ -C^T \Gamma^{-1} & \Sigma^{-1} + C^T \Gamma^{-1}C}\mat{z \\ w} -2\mat{z \\ w}^T\mat{\Gamma^{-1}d \\ \Sigma^{-1}\mu - C^T\Gamma^{-1}d} + \text{const}  \\
      &=\mat{z \\w}^T \mat{\Gamma + C \Sigma C^T & C\Sigma \\ \Sigma C^T & \Sigma}^{-1}\mat{z \\ w} -2\mat{z \\ w}^T\mat{\Gamma^{-1}d \\ \Sigma^{-1}\mu - C^T\Gamma^{-1}d} + \text{const}  \\
    \end{align*}
    Completing the square is:
    \[-\frac{1}{2} x\Sigma^{-1}x+b^Tx=-\frac{1}{2}(x-\Sigma b)^T\Sigma^{-1}(x-\Sigma b)+\text{const}
    \]
    So the mean is:
    \begin{align*}
    \mat{\mu_z \\ \mu_w} =\mat{\Gamma + C \Sigma C^T & C\Sigma \\ \Sigma C^T & \Sigma}\mat{\Gamma^{-1}d \\ \Sigma^{-1}\mu - C^T\Gamma^{-1}d}=\mat{C\mu+d \\ \mu}
    \end{align*}
    \item From (a) we can read off that:
    \[p(z|t) = \mathcal{N}(z|C \mu + d, \Gamma + C \Sigma C^T)= \mathcal{N}(z|C At + Cb + d, \Gamma + C \Sigma C^T)\]
    \item From (a) and wikipedia we deduce that
    $$p(w|z, t) = \Ncal(w|At+b + \Sigma C^T(\Gamma + C\Sigma C^T)^{-1}(z-C(At+b) - d), \Sigma - \Sigma C^T(\Gamma + C\Sigma C^T)^{-1}C\Sigma)$$
  \end{enumerate}
\end{gradedsolution}

\item \label{generating-sets}\textbf{BONUS}
Let $f: \mathcal{X} \to \mathcal{Y}$ be any function and $\mathcal{B}_{\mathcal{X}}$ a $\sigma$-algebra on $\mathcal{X}$.
Additionally, let $\mathcal{B}$ be any set of subsets of $\mathcal{Y}$, not necessarily a $\sigma$-algebra.
Assume that $f$ is $\mathcal{B}_{\mathcal{X}}$-$\mathcal{B}$-measurable, with which we mean that for each $B \in \mathcal{B}$, we have $f^{-1}(B) \in \mathcal{B}_{\mathcal{X}}$.
Show that $f$ is then already $\mathcal{B}_\mathcal{X}$-$\mathcal{B}_{\mathcal{Y}}$-measurable, where $\mathcal{B}_{\mathcal{Y}} := \sigma{(\mathcal{B})}$ is the smallest $\sigma$-algebra containing $\mathcal{B}$.

The bottom line, to be used in the second part of the exercise, is the following: measurability can be tested on sets which generate the $\sigma$-algebra of the codomain.

\emph{Hint: Show that $\mathcal{B} \subseteq f_*\mathcal{B}_{\mathcal{X}}$, where $f_*\mathcal{B}_{\mathcal{X}}$ is defined as in Definition/Lemma \ref*{sigma-algebra-map} in the appendix.
Deduce that $\mathcal{B}_{\mathcal{Y}} \subseteq f_* \mathcal{B}_{\mathcal{X}}$.
Deduce from that, in turn, that $f$ is $\mathcal{B}_{\mathcal{X}}$-$\mathcal{B}_{\mathcal{Y}}$-measurable.}

\begin{gradedsolution}
  Recall that
  \begin{equation*}
    f_* \mathcal{B}_{\mathcal{X}} = \{A \subseteq \mathcal{Y} \mid f^{-1}(A) \in \mathcal{B}_{\mathcal{X}}\}.
    \label{eq:recall_definition}
  \end{equation*}
  Since we have $f^{-1}(B) \in \mathcal{B}_{\mathcal{X}}$ for all $B \in \mathcal{B}$ by assumption, we obtain $\mathcal{B} \subseteq f_* \mathcal{B}_{\mathcal{X}}$.
  Now, since $f_* \mathcal{B}_{\mathcal{X}}$ is a $\sigma$-algebra containing $\mathcal{B}$, and since $\mathcal{B}_{\mathcal{Y}}$ is by definition the smallest $\sigma$-algebra containing $\mathcal{B}$, we obtain $\mathcal{B}_{\mathcal{Y}} \subseteq f_* \mathcal{B}_{\mathcal{X}}$.
  By definition of $f_* \mathcal{B}_{\mathcal{X}}$, again, this means that for all $B \in \mathcal{B}_{\mathcal{Y}}$ we have $f^{-1}(B) \in \mathcal{B}_{\mathcal{X}}$.
  This, in turn, means that $f$ is $\mathcal{B}_{\mathcal{X}}$-$\mathcal{B}_{\mathcal{Y}}$-measurable, proving the claim.
\end{gradedsolution}

\item \textbf{BONUS}
In this exercise we'll prove Lemma \ref*{transmeasure_is_twoargfunction} for Markov kernels/probability measures: Show that a Markov kernel, i.e.\ a measurable map
\[ K(W|T): \mathcal{T}\rightarrow \mathcal{P}(\mathcal{W}), \quad t\mapsto (D\mapsto K(W\in D|T=t)) \]
induces a two-argument function
\[ \tilde K(W|T):\, \Bcal_\Wcal \times \Tcal \to [0,1], \qquad (D,t) \mapsto \tilde K(W \in D|T=t) \]
such that:
\begin{enumerate}[label=\roman*)]
  \item for each $t \in \Tcal$ the map:
      \[ \Bcal_\Wcal \to [0,1],\qquad D \mapsto \tilde K(W\in D|T=t)\]
      is a probability measure, and
  \item for each $D \in \Bcal_\Wcal$ the map:
      \[ \tilde K(W \in D|T):\, \Tcal \to [0,1],\quad t \mapsto \tilde K(W \in D|T=t) \]
      is $\Bcal_\Tcal$-$\Bcal_{[0,1]}$-measurable
\end{enumerate}
and vice versa.

\emph{Hint: For doing this exercise, recall Definition \ref*{spaces_of_finite_and_prob_measures} of the $\sigma$-algebra of $\Pcal(\Wcal)$ and use exercise \ref*{generating-sets}.}
\begin{gradedsolution}
  \begin{enumerate}
    \item ($K\implies \tilde{K}$) Note that for each $t\in\Tcal$ we have $\tilde{K}(W | T=t) := K(W|T)(t) \in \mathcal{P}(\mathcal{W})$, which proves property i).
    Now, let $D \in \mathcal{B}_{\mathcal{W}}$ arbitrary and $t \in \mathcal{T}$.
    We obtain
    \begin{align*}
      \tilde{K}(W\in D \mid T=t) & = K(W|T)(t)(D) \\
      & = \text{ev}_D(K(t)) \\
      & = (\text{ev}_D \circ K)(t),
    \end{align*}
    which shows $\tilde{K}(W\in D \mid \bullet) = \text{ev}_D \circ K: \mathcal{T} \to [0,1]$.
    Since $K$ is measurable by assumption and $\text{ev}_D$ is measurable by the definition of the $\sigma$-algebra on $\mathcal{P}(\mathcal{W})$ given in Definition \ref*{spaces_of_finite_and_prob_measures}, and since compositions of measurable functions are measurable, we obtain the measurability of $\tilde{K}(W\in D \mid \bullet)$, which proves property ii).

    \item ($\tilde{K}\implies K$) Assume that $\tilde{K}: \mathcal{B}_{\mathcal{W}} \times \mathcal{T} \to [0,1]$ satisfies conditions i) and ii). 
    Define $K:\Tcal \rightarrow \Pcal(\Wcal): t\mapsto \tilde{K}(W|T=t)$. 
    If $t \in \mathcal{T}$, then $K(t) \in \mathcal{P}(\mathcal{W})$ according to condition i), which shows well-definedness of the map $K: \mathcal{T} \to \mathcal{P}(\mathcal{W})$.
    We need to show that it is also measurable, i.e.\ that preimages of measurable sets are measurable.
    Note that according to Definition \ref*{spaces_of_finite_and_prob_measures} the $\sigma$-algebra on $\mathcal{P}(\mathcal{W})$ is generated by sets of the form $\text{ev}_D^{-1}(A)$, where $D \subseteq \mathcal{W}$ is measurable, $A \subseteq [0,1]$ is measurable, and $\text{ev}_D: \mathcal{P(\mathcal{W})} \to [0,1]$ is an evaluation map.
    For these generating sets, we obtain
    \begin{align*}
      K^{-1}(\text{ev}_D^{-1}(A)) & = \{ t \in T \mid K(t) \in \text{ev}_D^{-1}(A)\} \\
      & = \{t \in T \mid \text{ev}_D(K(t)) \in A\} \\
      & = \{t \in T \mid K(t)(D) \in A\} \\
      & = \{t \in T \mid \tilde K(W\in D \mid T=t) \in A\} \\
      & = \tilde{K}(W\in D \mid \bullet)^{-1}(A) \in \Bcal_\Tcal
    \end{align*}
    according to condition ii). Using exercise \ref*{generating-sets} and that the sets $\text{ev}_D^{-1}(A)$ generate the $\sigma$-algebra on $\mathcal{P}(\mathcal{W})$, this proves the measurability of $K$.
  \end{enumerate}
\end{gradedsolution}

\end{enumerate}


\end{document}
